

The report-resolved continuation theorem is exact for linear marked task equalities of a supplied time-homogeneous finite-return routine. It does not identify the reachable positive cone with every positive matrix on its support, nor replace a one-sided reachable cost bound by an unsupported operator inequality. The reset-instrument classification applies to failure branches which physically prepare known states and to first-success commitment. It includes arbitrary finite rank-one families, but not arbitrary quantum instruments or an arbitrary-rank projection without input-erasing failure. The zero-risk spectrum is exact for the acquired profile; its physical threshold has the same converse and implementation hypotheses as the exact transfer theorem. These distinctions separate the new structural results from the independent wider research goals.
